\scanpage{82}
est un isomorphisme \inserted{en cohomologie} en bidegré $(i,j)$ avec $j\leq i$, tandis que la \unsure{coh.|ch.} du premier \unsure{mb} est nul pour $j>i$.

On peut l'exprimer comme un isomorphisme du $1^{\mathrm{er}}$ mb avec un certain tronqué du second....

\uline{Fonctorialités} a) en $k$, clair

b) \uline{En $u$}\deleted{\unsure{$\Phi$}} Soit
\[
u':\mathbb Z\longrightarrow\Phi'
\]
une deuxième immersion, S'il existe un hom $\lambda:\Phi\longrightarrow\Phi'$ rendant comm.
\diagram{\begin{tikzcd}[ampersand replacement=\&,column sep=small]
\& \mathbb Z\arrow[dl,"u"']\arrow[dr,"u'"] \& \& \\
\Phi\arrow[rr,"\lambda"] \& \& \Phi' \& \text{i.e.}\quad\lambda u=u'
\end{tikzcd}}
\deleted{\unsure{Eh}}On en conclut $\widetilde\lambda:K_u^{\bullet\bullet}\unsure{\mathord{-}|\to}K_{u'}^{\bullet\bullet}$ et $\widetilde\lambda_k:K_{u,k}^{\bullet\bullet}\longrightarrow K_{u',k}^{\bullet\bullet}$, compatible avec \unsure{ttes} les structures, qui est un isom. dans telle catégorie dérivée, et le reste après application de $\Gamma(X,-)$ ou de $f_*$.

Supposons qu'on ait un autre, soit $\mu$, tel que
\[
\mu u=u'
\]
alors, posant $\mu=\lambda+\Theta$,

on trouve $\Theta u=0$
% archive sheet number: 81
